\section{Niveau 2}

\begin{UPSTIinfor}{Une comparaison, ça vaut quelque chose}
    En C, une comparaison comme \texttt{x <= 10} n'est pas juste \og vraie \fg{} ou \og fausse \fg{} dans l'abstrait : elle est \textbf{évaluée} et produit la valeur entière \texttt{1} (vrai) ou \texttt{0} (faux). C'est cette valeur que le \texttt{if} regarde ensuite pour choisir sa branche.

    On peut aussi combiner une condition avec l'opérateur \textbf{NON} (\texttt{!} en C), qui inverse un résultat vrai/faux :
    \begin{center}
        \begin{tabular}{|c|c|}
            \hline
            \texttt{condition} & \texttt{NON condition} (\texttt{!condition}) \\
            \hline
            vrai                & faux                                          \\
            faux                & vrai                                          \\
            \hline
        \end{tabular}
    \end{center}
    Par exemple, \lstinline[language=C]{!(x > 10)} est vrai exactement quand \texttt{x > 10} est faux, c'est-à-dire quand \texttt{x <= 10}.
\end{UPSTIinfor}

\subsection{Conditions sur des intervalles}

\begin{UPSTIwarning}{Piège classique : enchaîner les comparaisons}
    En mathématiques, on écrit naturellement $2{,}0 \le p \le 5{,}0$. En C, \textbf{cette écriture ne fait pas ce que l'on croit} : \lstinline[language=C]{2.0 <= p <= 5.0} est évaluée de gauche à droite. On calcule d'abord \texttt{2.0 <= p}, qui vaut \texttt{0} ou \texttt{1} ; puis on compare \emph{ce résultat} (donc \texttt{0} ou \texttt{1}) à \texttt{5.0}. Comme \texttt{0} et \texttt{1} sont toujours inférieurs ou égaux à \texttt{5.0}, l'expression est \textbf{toujours vraie}, quelle que soit la valeur de \texttt{p} !
\end{UPSTIwarning}

\begin{UPSTIexercice}{Fenêtre de fonctionnement}
    Un capteur de pression n'est considéré comme fiable que si la pression mesurée \texttt{p} (en bar) respecte \(2{,}0 \le p \le 5{,}0\).
    \UPSTIquestion{Écrire, en langage C, la condition correcte à utiliser dans un \texttt{if} pour tester que le capteur est fiable.}
    \UPSTIquestion{Réécrire cette même condition d'une autre façon, en utilisant l'opérateur \texttt{!} et l'opérateur \texttt{||}.}
    \UPSTIquestion{Expliquer, en une phrase, pourquoi \lstinline[language=C]{p >= 2.0 <= 5.0} ne convient pas ici.}
\end{UPSTIexercice}

\begin{UPSTIprofOnlyEnv}
    \begin{UPSTIcorrectionP}{Fenêtre de fonctionnement}
        Condition directe : \lstinline[language=C]{(p >= 2.0 && p <= 5.0)}.

        Variante avec \texttt{!} et \texttt{||} : \lstinline[language=C]{!(p < 2.0 || p > 5.0)} (le capteur est fiable si l'on n'est \emph{ni} en dessous de 2.0 \emph{ni} au-dessus de 5.0).

        \lstinline[language=C]{p >= 2.0 <= 5.0} est évaluée de gauche à droite : \lstinline[language=C]{p >= 2.0} donne \texttt{0} ou \texttt{1}, puis ce résultat est comparé à \texttt{5.0} avec \texttt{<=}, ce qui est toujours vrai. Cette écriture ne teste donc jamais réellement la valeur de \texttt{p} par rapport à \texttt{5.0}.
    \end{UPSTIcorrectionP}
\end{UPSTIprofOnlyEnv}


\subsection{Combiner des conditions}

\begin{UPSTIexercice}{Évaluer des expressions booléennes}
    On considère trois entiers \texttt{a = 4}, \texttt{b = 9}, \texttt{c = 9}.
    \UPSTIquestion{Pour chaque expression ci-dessous, indiquer si elle vaut \textbf{vrai} ou \textbf{faux}.}
    \begin{multicols}{2}
        \begin{itemize}
            \item \verb|(a < b) && (b == c)|
            \item \verb$(a >= b) || (c != 9)$
            \item \verb|!(a < b)|
            \item \verb|!(a < b) && (b == c)|
            \item \verb$(a < b) || !(b == c)$
            \item \verb$!(a >= b) && !(c != 9)$
        \end{itemize}
    \end{multicols}
\end{UPSTIexercice}

\begin{UPSTIprofOnlyEnv}
    \begin{UPSTIcorrectionP}{Évaluer des expressions booléennes}
        Avec $a=4$, $b=9$, $c=9$ :
        \begin{itemize}
            \item \verb|(a < b) && (b == c)| : \texttt{vrai \&\& vrai} $\to$ \textbf{vrai}.
            \item \verb$(a >= b) || (c != 9)$ : \texttt{faux || faux} $\to$ \textbf{faux}.
            \item \verb|!(a < b)| : \texttt{a < b} est vrai, donc sa négation est \textbf{faux}.
            \item \verb|!(a < b) && (b == c)| : \texttt{faux \&\& vrai} $\to$ \textbf{faux}.
            \item \verb$(a < b) || !(b == c)$ : \texttt{vrai || faux} $\to$ \textbf{vrai}.
            \item \verb$!(a >= b) && !(c != 9)$ : \texttt{vrai \&\& vrai} $\to$ \textbf{vrai}.
        \end{itemize}
    \end{UPSTIcorrectionP}
\end{UPSTIprofOnlyEnv}


\begin{UPSTIexercice}{Ouverture d'une porte de laboratoire}
    Une porte de laboratoire s'ouvre automatiquement si la personne possède un \textbf{badge valide} (\texttt{bool badgeOK}) \textbf{et} a saisi un \textbf{code correct} (\texttt{bool codeOK}).
    \UPSTIquestion{Écrire, en pseudo-code, la condition qui décide de l'ouverture de la porte.}
    \UPSTIquestion{Traduire cette condition en C.}
    \UPSTIquestion{Un professeur peut aussi ouvrir la porte simplement en présentant son badge, même sans code (\texttt{bool estProfesseur}). Modifier la condition en C pour intégrer ce cas.}
\end{UPSTIexercice}

\begin{UPSTIprofOnlyEnv}
    \begin{UPSTIcorrectionP}{Ouverture d'une porte de laboratoire}
        Pseudo-code : \verb|badgeOK ET codeOK|.

        C : \lstinline[language=C]{(badgeOK && codeOK)}.

        Avec le cas du professeur : \lstinline[language=C]{((badgeOK && codeOK) || (estProfesseur && badgeOK))}, que l'on peut factoriser en \lstinline[language=C]{(badgeOK && (codeOK || estProfesseur))}.
    \end{UPSTIcorrectionP}
\end{UPSTIprofOnlyEnv}


\subsection{Modulo, division et conditions}

\begin{UPSTIexercice}{Multiples et conditions}
    Sur une chaîne de production, on veut déclencher un contrôle qualité approfondi toutes les 4 pièces produites. Le nombre de pièces déjà produites est stocké dans \texttt{int numPiece}.
    \UPSTIquestion{Que vaut \texttt{numPiece \% 4} lorsque \texttt{numPiece} est un multiple de 4 ? Et lorsqu'il ne l'est pas ?}
    \UPSTIquestion{Écrire, en C, la condition qui est vraie exactement lorsque \texttt{numPiece} est un multiple de 4.}
    \UPSTIquestion{Pour \texttt{numPiece} valant successivement 1, 2, 3, 4, 5, 6, 7, 8, indiquer pour lesquelles un contrôle est déclenché.}
    \UPSTIquestion{Écrire le code C complet qui affiche \og Controle qualite \fg{} lorsqu'un contrôle doit être déclenché, et qui ne fait rien sinon.}
\end{UPSTIexercice}

\begin{UPSTIprofOnlyEnv}
    \begin{UPSTIcorrectionP}{Multiples et conditions}
        \texttt{numPiece \% 4} vaut \texttt{0} lorsque \texttt{numPiece} est un multiple de 4, et vaut \texttt{1}, \texttt{2} ou \texttt{3} sinon.

        Condition : \lstinline[language=C]{(numPiece % 4 == 0)}.

        Un contrôle est déclenché pour \texttt{numPiece = 4} et \texttt{numPiece = 8} uniquement (les multiples de 4 dans la liste).

        \begin{lstlisting}[language=C]
if (numPiece % 4 == 0)
{
    printf("Controle qualite\n");
}
        \end{lstlisting}
    \end{UPSTIcorrectionP}
\end{UPSTIprofOnlyEnv}

\begin{UPSTIwarning}{Le modulo ne s'utilise qu'avec des entiers}
    L'opérateur \texttt{\%} n'est défini, en C, qu'entre deux entiers. Écrire \lstinline[language=C]{p % 1.0}, où \texttt{p} est un \texttt{float}, ne compile pas : le compilateur refuse un \texttt{\%} sur des valeurs flottantes. Pour tester si une mesure flottante (une pression, une tension...) respecte un pas régulier, il faut donc passer par autre chose qu'un simple \texttt{\%} --- par exemple ramener le problème à des entiers, ou comparer un écart à une tolérance.
\end{UPSTIwarning}


\subsection{Enchaînements à plusieurs branches}

\begin{UPSTIexercice}{Classer une résistance par bande de valeurs}
    Un programme classe une résistance mesurée \texttt{r} (en ohms) :
    \begin{lstlisting}[language=C]
if (r < 100)
{
    printf("Faible\n");
}
else if (r < 1000)
{
    printf("Moyenne\n");
}
else if (r < 10000)
{
    printf("Elevee\n");
}
else
{
    printf("Tres elevee\n");
}
    \end{lstlisting}
    \UPSTIquestion{Pour \texttt{r = 850}, indiquer, dans l'ordre, toutes les conditions réellement testées avant l'affichage.}
    \UPSTIquestion{Combien de conditions au maximum peuvent être testées pour une valeur de \texttt{r}, quelle qu'elle soit ? Justifier.}
    \UPSTIquestion{On intervertit les deux premiers blocs \texttt{if} et \texttt{else if} (celui de \og Faible \fg{} et celui de \og Moyenne \fg), \textbf{sans changer les conditions}. Le programme obtenu classe-t-il encore correctement \texttt{r = 850} et \texttt{r = 50} ? Justifier.}
\end{UPSTIexercice}

\begin{UPSTIprofOnlyEnv}
    \begin{UPSTIcorrectionP}{Classer une résistance par bande de valeurs}
        Pour \texttt{r = 850} : on teste \texttt{r < 100} (faux), puis \texttt{r < 1000} (vrai) $\to$ affichage de \og Moyenne \fg. Seules ces deux conditions sont testées.

        Au maximum, les trois conditions peuvent être testées (si aucune des trois premières n'est vraie, on tombe dans le \texttt{else} final sans test supplémentaire).

        En intervertissant les deux premiers blocs, on obtiendrait :
        \begin{lstlisting}[language=C]
if (r < 1000)      { printf("Moyenne\n"); }
else if (r < 100)  { printf("Faible\n"); }
        \end{lstlisting}
        Pour \texttt{r = 850} : \texttt{r < 1000} est vrai, on affiche toujours \og Moyenne \fg{} $\to$ correct. Mais pour \texttt{r = 50} : \texttt{r < 1000} est déjà vrai (50 < 1000), donc on affiche \og Moyenne \fg{} au lieu de \og Faible \fg{} $\to$ \textbf{incorrect}. Le second \texttt{else if} ne sera jamais atteint pour une petite valeur, car la première condition, trop large, l'a déjà \og capturée \fg. \textbf{L'ordre des conditions dans un enchaînement compte.}
    \end{UPSTIcorrectionP}
\end{UPSTIprofOnlyEnv}


\begin{UPSTIexercice}{Traduire un enchaînement complet}
    Un système de ventilation affiche un message selon la température \texttt{temp} et l'humidité \texttt{humid} :
    \begin{itemize}
        \item si \texttt{temp > 30} \textbf{ET} \texttt{humid > 70}, afficher \og Danger canicule humide \fg{} ;
        \item sinon, si \texttt{temp > 30} (mais \texttt{humid <= 70}), afficher \og Chaleur seche \fg{} ;
        \item sinon, afficher \og Conditions normales \fg.
    \end{itemize}
    \UPSTIquestion{Écrire le pseudo-code complet de ce système, avec \texttt{SI}, \texttt{SINON SI}, \texttt{SINON}.}

    \UPSTIquestion{Traduire ce pseudo-code en C.}

\end{UPSTIexercice}

\begin{UPSTIprofOnlyEnv}
    \begin{UPSTIcorrectionP}{Traduire un enchaînement complet}
        \begin{verbatim}
SI temp > 30 ET humid > 70 ALORS
    AFFICHER "Danger canicule humide"
SINON SI temp > 30 ALORS
    AFFICHER "Chaleur seche"
SINON
    AFFICHER "Conditions normales"
FIN SI
        \end{verbatim}
        \begin{lstlisting}[language=C]
if (temp > 30 && humid > 70)
{
    printf("Danger canicule humide\n");
}
else if (temp > 30)
{
    printf("Chaleur seche\n");
}
else
{
    printf("Conditions normales\n");
}
        \end{lstlisting}
        Remarque : dans le second bloc, il est inutile de retester \texttt{humid <= 70} : on n'y arrive que si la première condition (qui contenait déjà \texttt{temp > 30}) était fausse, et comme \texttt{temp > 30} vient d'être confirmé vrai ici, c'est forcément \texttt{humid > 70} qui était fausse.
    \end{UPSTIcorrectionP}
\end{UPSTIprofOnlyEnv}
